\section*{Chapitre \ref{ch:g11:scal} --- Le produit scalaire dans le plan}

\begin{solution}{exo:g11:scal:1}
$(2,5)\cdot(3,-1) = 6 - 5 = 1$.

$(1,-4)\cdot(8,2) = 8 - 8 = 0$ (\omterm{def:g11:scal:orthogonal}{vecteurs orthogonaux}).

$\vec u\cdot\vec v = 3 \times 4 \times \cos\frac\pi3 = 12 \times \frac12
= 6$.
\end{solution}

\begin{solution}{exo:g11:scal:2}
$\vect{AB}\cdot\vect{AC} = 6 \times 6 \times \cos 60^\circ = 18$.

Les \omterm{def:g11:vect:vector}{vecteurs} $\vect{AB}$ et $\vect{BC}$ font un angle de
$180^\circ - 60^\circ = 120^\circ$ (les placer bout à bout en $B$), donc
$\vect{AB}\cdot\vect{BC} = 36 \cos 120^\circ = -18$.
\end{solution}

\begin{solution}{exo:g11:scal:3}
$\vec u \cdot \vec v = 3(t - 8) + t = 4t - 24 = 0$ donne $t = 6$.

$\vec u - \vec v = (t - 4,\ 0)$, donc
$\vec u \cdot (\vec u - \vec v) = t(t-4) + 0 = t^2 - 4t$, qui s'annule
pour $t = 0$ ou $t = 4$.
\end{solution}

\begin{solution}{exo:g11:scal:4}
$\vec u \cdot \vec v = 3 + 2 = 5$, $\norm{\vec u} = \sqrt5$,
$\norm{\vec v} = \sqrt{10}$, donc
\[
\cos\theta = \frac{5}{\sqrt5\,\sqrt{10}} = \frac{5}{5\sqrt2}
= \frac{\sqrt2}{2},
\]
d'où $\theta = 45^\circ$ exactement.
\end{solution}

\begin{solution}{exo:g11:scal:5}
$a^2 = 16 + 36 - 2 \times 4 \times 6 \times \cos\frac{2\pi}{3}
= 52 - 48 \times \left(-\frac12\right) = 76$, donc $a = \sqrt{76} =
2\sqrt{19} \approx 8.7$.

Pour le second triangle, la loi des cosinus résolue pour l'angle donne
\[
\cos\widehat A = \frac{b^2 + c^2 - a^2}{2bc}
= \frac{25 + 16 - 49}{2 \times 5 \times 4} = -\frac{8}{40} = -\frac15,
\]
négatif~: $\widehat A$ est obtus.
\end{solution}

\begin{solution}{exo:g11:scal:6}
$\vect{AB}\,(4, -2)$ et $\vect{AC}\,(2, 4)$~:
$\vect{AB}\cdot\vect{AC} = 8 - 8 = 0$, donc l'angle en $A$ est droit.
Côtés~: $AB = \sqrt{16 + 4} = \sqrt{20}$, $AC = \sqrt{4 + 16} =
\sqrt{20}$, et $\vect{BC}\,(-2, 6)$ donne $BC = \sqrt{40}$.
Cohérence~: $BC^2 = 40 = 20 + 20 = AB^2 + AC^2$, la loi des cosinus avec
$\cos\widehat A = 0$ (Pythagore).
\end{solution}

\begin{solution}{exo:g11:scal:7}
Forme centre--rayon~: $(x + 2)^2 + (y - 3)^2 = 25$.

Pour le second cercle, compléter les carrés~:
\[
x^2 - 6x + y^2 + 2y = 6
\iff (x - 3)^2 - 9 + (y + 1)^2 - 1 = 6
\iff (x-3)^2 + (y+1)^2 = 16 :
\]
centre $(3, -1)$, rayon $4$.
\end{solution}

\begin{solution}{exo:g11:scal:8}
$M(x,y)$ est sur le cercle lorsque $\vect{MA}\cdot\vect{MB} = 0$ avec
$\vect{MA}\,(-3 - x,\ -y)$ et $\vect{MB}\,(-x,\ 2 - y)$~:
\[
(-3 - x)(-x) + (-y)(2 - y) = x^2 + 3x + y^2 - 2y = 0 .
\]
Pour l'origine~: $0 + 0 + 0 - 0 = 0$, donc $O$ est sur le cercle ---
géométriquement, $\vect{OA}\,(-3,0)$ et $\vect{OB}\,(0,2)$ sont
\omterm{def:g11:scal:orthogonal}{orthogonaux}, donc $O$ voit $[AB]$ sous un angle droit. (En complétant les
carrés~: centre $\left(-\frac32, 1\right)$, rayon
$\frac{\sqrt{13}}{2}$.)
\end{solution}

\begin{solution}{exo:g11:scal:9}
La perpendiculaire par $P$ a le \omterm{def:g11:vect:direction}{vecteur directeur} de $d$, à savoir
$(-1, 2)$, comme \omterm{def:g11:vect:vector}{vecteur} normal~: \omterm{def:g10:algebra:equation}{équation} $-x + 2y + c = 0$~; par
$P(1,3)$~: $-1 + 6 + c = 0$, $c = -5$, donc $x - 2y + 5 = 0$.

\omterm{def:g10:numbers:interunion}{Intersection} avec $d$~: de $2x + y = 7$, $y = 7 - 2x$~; en substituant,
$x - 2(7 - 2x) + 5 = 5x - 9 = 0$, donc $x = \frac95$, $y = \frac{17}5$~:
$H\left(\frac95, \frac{17}5\right)$. Alors
$\vect{PH}\left(\frac45, \frac25\right)$ et
\[
PH = \sqrt{\frac{16}{25} + \frac{4}{25}} = \frac{\sqrt{20}}{5}
= \frac{2\sqrt5}{5} \approx 0.89 .
\]
\end{solution}

\begin{solution}{exo:g11:scal:10}
En additionnant
\[
\norm{\vec u + \vec v}^2 = \norm{\vec u}^2 + 2\vec u\cdot\vec v +
\norm{\vec v}^2
\quad\text{et}\quad
\norm{\vec u - \vec v}^2 = \norm{\vec u}^2 - 2\vec u\cdot\vec v +
\norm{\vec v}^2,
\]
les termes croisés s'annulent et
$\norm{\vec u + \vec v}^2 + \norm{\vec u - \vec v}^2 = 2\norm{\vec u}^2 +
2\norm{\vec v}^2$. Dans un parallélogramme de côtés $\vec u$ et $\vec v$,
les diagonales sont $\vec u + \vec v$ et $\vec u - \vec v$~: la somme des
carrés des deux diagonales égale la somme des carrés des quatre côtés.
\end{solution}

\begin{solution}{exo:g11:scal:11}
\emph{1.} Comme $I$ est le \omterm{prop:g10:coordgeom:midpoint}{milieu}, $\vect{IB} = -\vect{IA}$ et
$IA = 2$. Alors
\[
\vect{MA}\cdot\vect{MB}
= (\vect{MI} + \vect{IA})\cdot(\vect{MI} - \vect{IA})
= \norm{\vect{MI}}^2 - \norm{\vect{IA}}^2
= MI^2 - 4 .
\]

\emph{2.} La condition $\vect{MA}\cdot\vect{MB} = 12$ devient
$MI^2 = 16$, c'est-à-dire $MI = 4$~: l'ensemble est le cercle de centre
$I$ et de rayon $4$.
\end{solution}

\begin{solution}{pb:g11:scal:1}
\textbf{1.} $\vec u \cdot \vec v = -3 + 8 = 5$~;
$\norm{\vec u} = 5$, $\norm{\vec v} = \sqrt5$~;
$\cos\theta = \frac{5}{5\sqrt5} = \frac{1}{\sqrt5}$, d'où
$\theta \approx 63.4^\circ$.

\textbf{2.} $6 - 4k = 0$~: $k = \frac32$.

\textbf{3.} $W = 50 \times 10 \times \cos 60^\circ = 250$ joules~:
seule la composante le long du déplacement travaille.

\textbf{4.} $c^2 = 25 + 49 - 2 \times 35 \times \frac12 = 39$~:
$c = \sqrt{39}$.

\textbf{5.} $\cos\gamma = \frac{16 + 36 - 64}{2 \times 4
\times 6} = -\frac14$~: négatif, l'angle est obtus~:
$\gamma \approx 104.5^\circ$.

\textbf{6.} Par les \omterm{def:g10:coordgeom:system}{coordonnées}~: $\vec u \cdot \vec v = \cos a
\cos b + \sin a \sin b$. Par la géométrie~: deux \omterm{def:g11:vect:vector}{vecteurs}
unitaires, donc $\vec u \cdot \vec v = 1 \times 1 \times
\cos(a - b)$. En égalant les deux calculs d'un même nombre~:
$\cos(a - b) = \cos a \cos b + \sin a \sin b$.

\textbf{7.} $\cos(a + b) = \cos(a - (-b)) = \cos a \cos(-b) +
\sin a \sin(-b) = \cos a \cos b - \sin a \sin b$.

\textbf{8.} $\sin(a + b) = \cos\left(\frac\pi2 - a - b\right)
= \cos\left(\left(\frac\pi2 - a\right) - b\right)
= \cos\left(\frac\pi2 - a\right)\cos b +
\sin\left(\frac\pi2 - a\right)\sin b
= \sin a \cos b + \cos a \sin b$.

\textbf{9.} $\cos 15^\circ = \cos 45^\circ \cos 30^\circ +
\sin 45^\circ \sin 30^\circ = \frac{\sqrt2}{2} \cdot
\frac{\sqrt3}{2} + \frac{\sqrt2}{2} \cdot \frac12 =
\frac{\sqrt6 + \sqrt2}{4} \approx 0.966$. Et
$\sin 75^\circ = \cos 15^\circ$~: même valeur.

\textbf{10.} En posant $b = a$ dans les questions 7 et 8~:
$\cos 2a = \cos^2 a - \sin^2 a = 2\cos^2 a - 1 = 1 - 2\sin^2 a$
(à l'aide de $\cos^2 + \sin^2 = 1$), et
$\sin 2a = 2 \sin a \cos a$.

\textbf{11.} $\sin^2 15^\circ = \frac{1 - \cos 30^\circ}{2}
= \frac{2 - \sqrt3}{4}$, donc
$\sin 15^\circ = \frac{\sqrt{2 - \sqrt3}}{2}$. En élevant
$\frac{\sqrt6 - \sqrt2}{4}$ au carré~:
$\frac{6 - 2\sqrt{12} + 2}{16} = \frac{8 - 4\sqrt3}{16} =
\frac{2 - \sqrt3}{4}$~: le même nombre en deux costumes.

\textbf{12.} $\vect{AP}$ se décompose en une part le long de $d$
(invisible pour $\vec n$) et une part perpendiculaire à $d$, dont la
longueur est la distance cherchée~; multiplier scalairement par la
normale unitaire $\frac{\vec n}{\norm{\vec n}}$ annule la première
part et mesure la seconde. Avec $A$ sur $d$~:
$\vect{AP} \cdot \vec n = a x_0 + b y_0 - (a x_A + b y_A) =
a x_0 + b y_0 + c$, d'où la formule. Exemple~:
$\frac{\abs{15 + 24 - 12}}{5} = \frac{27}{5} = 5.4$.

\textbf{13.} Rayon $= 5.4$~: $(x - 5)^2 + (y - 6)^2 = 29.16$.

\textbf{14.} Distance $= \frac{\abs{-12}}{5} = 2.4$. Pied de la
perpendiculaire~: la perpendiculaire issue de l'origine est
$(3t, 4t)$~; sur la droite~: $9t + 16t = 12$, donc
$t = \frac{12}{25}$, et le pied est
$\left(\frac{36}{25}, \frac{48}{25}\right)$, à la distance
$5t = \frac{60}{25} = 2.4$. La formule et le calcul honnête
concordent.

\textbf{15.} Base $AB = 6$, hauteur $=$ distance de $C$ à $(AB)$
(l'axe $y = 0$), soit $5$~; aire $15$. Droite $(AC)$~:
$5x - 2y = 0$, et
$\operatorname{dist}(B, (AC)) = \frac{\abs{30}}{\sqrt{29}}$~;
alors $\frac12 \times AC \times \frac{30}{\sqrt{29}} =
\frac12 \times \sqrt{29} \times \frac{30}{\sqrt{29}} = 15$~: la
même aire, quelle que soit la hauteur choisie.

\textbf{16.} $\operatorname{dist} = \frac{\abs{2 + 1 - 6}}
{\sqrt2} = \frac{3}{\sqrt2} \approx 2.12 > 2$~: la droite passe au
large du cercle --- ils sont disjoints.

\textbf{17.} En insérant $I$~:
$\vect{MA} \cdot \vect{MB} = (\vect{MI} + \vect{IA}) \cdot
(\vect{MI} + \vect{IB}) = MI^2 + \vect{MI} \cdot (\vect{IA} +
\vect{IB}) + \vect{IA} \cdot \vect{IB}$. Le terme du \omterm{prop:g10:coordgeom:midpoint}{milieu}
disparaît ($\vect{IA} + \vect{IB} = \vec 0$), et
$\vect{IA} \cdot \vect{IB} = -IA^2 = -\frac{AB^2}{4}$ (\omterm{def:g11:vect:vector}{vecteurs}
opposés de longueur $\frac{AB}2$).

\textbf{18.} $\vect{MA} \cdot \vect{MB} = 0 \iff MI^2 =
\frac{AB^2}{4} \iff MI = \frac{AB}{2}$~: c'est le cercle de centre
$I$ et de rayon $\frac{AB}2$ --- le cercle de diamètre $[AB]$.

\textbf{19.} En développant chaque couple à partir de $H$
($\vect{BC} = \vect{HC} - \vect{HB}$, etc.)~:
\[
\vect{HA} \cdot \vect{HC} - \vect{HA} \cdot \vect{HB}
+ \vect{HB} \cdot \vect{HA} - \vect{HB} \cdot \vect{HC}
+ \vect{HC} \cdot \vect{HB} - \vect{HC} \cdot \vect{HA} = 0 :
\]
tout s'annule deux à deux --- l'identité vaut pour \emph{quatre}
points quelconques. Pour notre $H$, les deux premiers \omterm{def:g11:scal:dot}{produits
scalaires} sont nuls par hypothèse, donc
$\vect{HC} \cdot \vect{AB} = 0$~: $H$ appartient à la hauteur issue
de $C$. Trois hauteurs, un seul point.

\textbf{20.} Les formules d'addition sont venues du calcul d'un même
\omterm{def:g11:scal:dot}{produit scalaire} par le visage des \emph{\omterm{def:g10:coordgeom:system}{coordonnées}} et par celui
des \emph{\omterm{def:g11:scal:dot}{normes} et de l'angle}~; la formule de distance, du visage
de la \emph{projection}~; l'\omterm{pb:g11:scal:1}{orthocentre}, du visage de la
\emph{bilinéarité} (\omterm{def:g10:algebra:expand}{développer} avec Chasles). Calculer deux fois la
même quantité, égaler, et un théorème en tombe --- c'est tout le
métier du \omterm{def:g11:scal:dot}{produit scalaire}.
\end{solution}
